% Generated by report_gen.py (ai-ee v1) - do not hand-edit.
\documentclass[11pt,a4paper]{article}
\usepackage[margin=2.2cm]{geometry}
\usepackage{graphicx}
\usepackage{booktabs}
\usepackage{longtable}
\usepackage{pdfpages}
\makeatletter
\@ifundefined{KV@Gin@artifact}{\define@key{Gin}{artifact}[]{}}{}
\makeatother
\usepackage{xcolor}
\usepackage[hidelinks]{hyperref}
\setcounter{tocdepth}{1}
\begin{document}
\sloppy
\begin{center}
{\LARGE\bfseries bb-adc --- Design Document}\\[6pt]
{\large ai-ee v1 pipeline}\\[2pt]
generated 2026-08-16 21:31:16
\end{center}
\tableofcontents

\section{Board and Run Metadata}
{\small
\begin{longtable}{lp{11cm}}
\toprule
\textbf{Field} & \textbf{Value} \\
\midrule
\endhead
board & bb-adc \\
workspace & \texttt{boards/bb-adc} \\
phase & P2 \\
gate status & no gates recorded yet \\
state created & 2026-08-16T17:25:12 \\
state updated & 2026-08-16T21:25:47 \\
\bottomrule
\end{longtable}
}
\section{Overview}
learning block-basics: a single-channel analog-to-digital converter. Analog input 0-5 V DC on a screw terminal; 12-bit or better; up to 10 kSa/s; DC accuracy matters more than speed. Digital output to a host over a 0.1 inch header, which also feeds the board its 3.3 V rail. Workspace boards/bb-adc.

---

Owner brief, verbatim, received 2026-08-16. The text above is the owner's own wording and is the authoritative statement of the task; nothing below it is part of the brief.

Mode token resolved at P0 by \texttt{state.py mode}:

\begin{itemize}
\item token: \texttt{learning block-basics:}
\item target: \texttt{block-basics} -\textgreater{} scope tier \texttt{block-only}, binding \texttt{canonical}
\item geometry: OUTPUT (the placement earns the outline; \texttt{board\_init --outline auto})
\item stage under study: none
\end{itemize}
\subsection*{Stackup}
Chosen stackup: \texttt{JLC2313\_1.6}
\section{Requirements}
\subsection*{Requirements: bb-adc}
\subsubsection*{1. Function}
A single-channel analog-to-digital converter, built as a study article. It takes one DC analog voltage in the range 0 to 5 V on a screw terminal, digitizes it to 12 bits or better at a rate up to 10 kSa/s, and hands the result to a host over a 0.1 inch header. That same header feeds the board its only supply, 3.3 V. DC accuracy is the dominant requirement; conversion speed is secondary and explicitly subordinate to it. Nothing else is on the board: no MCU, no display, no indicators, no second channel.

BUILD MODE: the brief opens with the token \texttt{learning block-basics:}. Resolved at P0 (\texttt{state.py mode}, recorded as a P0 decision; contract in \texttt{reference/build-modes.md}):

\begin{itemize}
\item mode / target: \textbf{block-basics}
\item scope tier: \textbf{block-only}
\item binding level: \textbf{canonical} - board geometry is an \textbf{OUTPUT} of the design,
not an input to it
\item stage under study: \textbf{none} (the block end to end, no single stage frozen as
a bench fixture)
\end{itemize}
What that means here. The block under study is the ADC signal chain itself: the converter, exactly the support parts its datasheet requires for correct operation at the stated operating point, one input interface (the screw terminal) and one output interface (the host header), plus what the fab needs to build the board and the bench needs to hold it. Protection of every kind, filtering the datasheet does not require, indicators, config straps, test points beyond the block's own measurement need, any second rail or second IC the block does not need in order to work, and mechanical/enclosure features beyond mounting the bare board are OUT by MODE, not by engineering judgment - and must never be recorded as generally unnecessary. The mode's defaults are applied below (quantity 5, JLC PCBA single-sided, bench 0-50 C indoor, no enclosure, fewest honest layers, geometry earned by the placement, stop at P9) instead of being asked. The mode relaxes geometry, cost and packaging only: every electrical spec, every gate, the coverage checks and the research they trigger, DFM, the datasheet's own requirements and every safety question in section 8 stand unchanged.

\textbf{Recorded requirement tension, NOT solved here (P1 owns it).} The stated analog input span (0 to 5 V) sits ABOVE the board's only stated rail (3.3 V from the host header). Every way out of that is an architecture decision - which converter, which reference, whether the input is attenuated ahead of the converter, whether the block honestly needs something the 3.3 V rail cannot give it - and the architect decides it, not this document. Two consequences that must travel with the tension:

\begin{itemize}
\item It touches a MODE BOUNDARY. \texttt{block-only} excludes "any second rail the block
does not need in order to work"; if P1 concludes the block DOES need one to
work at the stated operating point, that is support, not a feature, and the
call belongs at H1 with the owner, recorded as a decision.
\item It interacts with the accuracy answer (Open question 2). Whatever sits
between the terminal and the converter is inside the DC error budget, so the
topology choice and the accuracy number have to be settled together.
\end{itemize}
\subsubsection*{2. Interfaces}
Exactly two external interfaces, which is what the scope tier allows and what the brief states.

\begin{itemize}
\item \textbf{J1 - analog input.} Screw terminal (STATED). ASSUMED: 2-pole through-hole
terminal block, 5.08 mm pitch - signal and its ground return - rated far
above anything this board carries; pitch was not stated and 5.08 mm is the
common, well-stocked JLC/LCSC choice that accepts bench lead wire. ASSUMED:
single-ended input referenced to board ground, since the brief describes one
channel of 0-5 V DC with no mention of a floating or differential source
(confirm via Open question 4). Silkscreen-marked for signal and ground.
NOTE: the board has no input protection by mode, so the silk marking and the
owner's answer to Open question 5 are the only defenses against a wrong
connection - see section 8.
\item \textbf{J2 - host header.} 0.1 inch (2.54 mm) pitch header (STATED). It carries
BOTH the digital output and the board's 3.3 V supply plus ground (STATED).
ASSUMED: single-row, through-hole male pin header, one row, on a board edge.
Pin count, pin order and protocol are NOT stated - see Open questions 3
and 6.
\item \textbf{Digital output signalling.} Protocol NOT stated (Open question 3).
ASSUMED once the protocol is known: 3.3 V CMOS levels, referenced to the
same ground the header supplies, host acts as the bus master and the board
is a slave/peripheral - it never drives the interface unasked. ASSUMED:
conversions are host-triggered (the host asks, the board answers); a
dedicated data-ready or convert-start pin is added only if the chosen
converter genuinely requires one, in which case it lands on J2 and the pin
count in Open question 6 grows by that one line.
\item \textbf{Sample rate.} ASSUMED: "up to 10 kSa/s" is a MAXIMUM capability, not a
sustained obligation. Any rate from single-shot on demand up to 10 kSa/s
must work; nothing requires gap-free streaming at the top rate. Low risk to
assume, since a lower rate never makes a DC-accurate converter harder.
\item \textbf{No other external interface exists or is wanted}: no USB, no second
channel, no LED, no button, no jumper, no config strap, no external
reference input connector, no test points beyond the block's own measurement
need. All excluded by mode.
\end{itemize}
\subsubsection*{3. Power}
\begin{itemize}
\item \textbf{Only rail: 3.3 V, supplied BY the host through J2 (STATED).} The board
generates no rail of its own by default. Nominal 3.3 V; tolerance NOT stated
\begin{itemize}
\item ASSUMED 3.3 V +/-5 \% (3.135 to 3.465 V) at the header pin, which is the
ordinary tolerance of a regulated host rail. Confirm with Open question 7.
\end{itemize}
\item \textbf{Rail quality is an accuracy input, not a detail.} If the design ends up
referencing the converter to its own supply, the host's rail noise and
tolerance land directly in the measurement error. Whether that is acceptable
depends on the answer to Open question 2 and is a P1/P2 decision; it is
flagged here so the choice is deliberate.
\item \textbf{Current budget (GUESS, for sizing and for Open question 7 only).}
Converter 0.5 to 5 mA depending on class and rate; a precision voltage
reference, if the accuracy answer requires one, 0.1 to 1.5 mA; an input
buffer, if the source-impedance answer requires one, 0.5 to 2 mA; leakage
and bias elsewhere well under 1 mA. Total GUESS: under 10 mA typical, under
15 mA worst case. Marked as a guess: no part is chosen yet.
\item {*}{*}No battery, no charging, no energy storage, no backup rail, no second
rail{*}{*} on the board as specified. A second rail is excluded by the scope
tier unless P1 shows the block cannot work without one - see the tension
recorded in section 1.
\item \textbf{Board dissipation}: negligible (GUESS: under 50 mW). This board has no
thermal problem to solve. Temperature still matters, but as a DRIFT term in
the accuracy budget over the 0-50 C ambient (section 4), not as a cooling
requirement.
\item {*}{*}Excluded by mode (SCOPE decisions - do not record any of these as
generally unnecessary){*}{*}: input TVS/ESD clamp, series/reverse protection,
fusing, local rail filtering beyond the decoupling the converter's datasheet
requires, and any supply sequencing or supervisor.
\end{itemize}
\subsubsection*{4. Environment}
Mode defaults - applied, not asked.

\begin{itemize}
\item Ambient: 0 to 50 C, indoor bench.
\item Cooling: natural convection only. No heatsink, no forced airflow, no chassis
thermal path. Nothing on this board needs one.
\item No enclosure. No ingress (IP) rating. No vibration or shock requirement. No
humidity requirement. No conformal coating.
\item No formal EMC/emissions campaign - this is a bench study article, not a
product. Clean analog return geometry is still required because it is what
makes the block correct, not because a test demands it.
\item \textbf{Consequence that binds the electrical spec:} the 0-50 C range is what
turns reference tempco, amplifier offset drift and (if an attenuator is
used) resistor TCR mismatch into real, budgetable error terms. The accuracy
answer in Open question 2 is therefore stated at 25 C AND across 0-50 C.
\end{itemize}
\subsubsection*{5. Size \& mounting}
\begin{itemize}
\item \textbf{Outline: no HARD cap. RELAXABLE (canonical).} The binding level is
\texttt{canonical}, so board size, aspect and outline are OUTPUTS of the placement.
Nothing in this section binds at P5 \texttt{board\_init}. No planning dimension is
written here on purpose: bb-buck recorded a "for planning only" figure, a
number reached the CLI two checkpoints later, and placement optimized to fit
it - the board ended up with 0.05 mm of slack on all four edges and taught
the wrong lesson. This board will not repeat that.
\item \textbf{The flow that earns the outline} (\texttt{reference/build-modes.md}, binding
levels):
1. P5 \texttt{board\_init --outline auto} - generous provisional room. A fixed
\texttt{--outline WxH} is REFUSED under this mode, and that refusal is correct.
2. P6 place to the canonical layout, gate \texttt{place}.
3. P6 close: \texttt{board\_edit --outline fit --margin M} - the board becomes what
the placement needs. Re-run \texttt{planes\_gen} if it GREW.
4. P7 route.
\item \textbf{Layers}: ASSUMED 2, per the mode's "fewest layers the block honestly
needs". If P2/P6 shows that an honest analog return and reference layout
cannot be had on 2 layers, 4 is allowed under the same rule. Not fixed here
either way.
\item \textbf{Mounting}: ASSUMED 4 x M3 clearance holes (3.2 mm) inset from the
corners, so the board can sit on standoffs at the bench. Mounting the bare
board is explicitly in scope; enclosure fit is not. If the earned outline is
too small for four, two on opposite corners is acceptable.
\item \textbf{Height}: ASSUMED no maximum. The board is open on a bench with no
enclosure, and the screw terminal itself stands roughly 10-12 mm tall with
wire fitted, so a height cap would be arbitrary.
\item \textbf{Connector placement}: ASSUMED the screw terminal and the host header sit
on different board edges with their openings facing outward, so the analog
input wiring does not have to cross the digital header to reach the bench.
Which edges, and the geometry inside the board, is the layout engineer's
call at P6 - that is the whole point of a canonical binding.
\end{itemize}
\subsubsection*{6. Quantity \& budget}
Mode defaults - applied, not asked.

\begin{itemize}
\item Build quantity: 5.
\item Cost: minimal, with no hard per-unit cap. Prefer JLC Basic/Preferred parts
wherever a Basic part actually meets the electrical requirement; Extended is
allowed for the converter and for a precision voltage reference, where Basic
stock in the required accuracy class is thin. Cost is one of the three
things this mode may relax, and accuracy is not - if the accuracy answer and
a Basic part conflict, accuracy wins and the choice is recorded.
\end{itemize}
\subsubsection*{7. Assembly}
Mode defaults - applied, not asked.

\begin{itemize}
\item JLCPCB PCBA, single-sided: all SMT parts on the top side only. The bottom
side stays clear of SMT so it can serve as unbroken return copper.
\item The screw terminal and the 0.1 inch header are through-hole. ASSUMED: fitted
by JLCPCB through-hole assembly if offered at acceptable cost for the chosen
parts, hand-soldered after SMT otherwise. Either path leaves the schematic
and the footprints unchanged.
\item Pipeline stops at P9 (fab package + DFM). Ordering is a separate owner
decision and is not requested here.
\end{itemize}
\subsubsection*{8. Compliance/safety flags}
The mode grants no silence in this section. One safety-relevant fact is genuinely missing from the brief and is asked as Open question 1; it is not guessed, and the pipeline should not proceed past H1 without it.

\begin{itemize}
\item \textbf{Mains voltage: DOES NOT APPLY to anything ON this board} - the highest
voltage on any board net at the stated operating point is 5 V. {*}{*}BUT the
source of the 0-5 V input is NOT stated{*}{*}, and that is the one thing that
could make this section wrong. A 0-5 V signal is very often the low side of
something much larger: a shunt amplifier in a mains-referenced circuit, a
divider off a high-voltage bus, an isolated sensor's output that turns out
not to be isolated. This board has no isolation and no input protection by
mode, and its input ground is the HOST's ground through J2, so a hazardous
or merely elevated source ties the host to that source directly. ASKED, not
assumed - Open question 1.
\item \textbf{Batteries: DOES NOT APPLY.} No battery, no charging function, no pack, no
cell chemistry, no deep-discharge behavior anywhere on this board. The board
is fed by the host's 3.3 V rail, which it does not generate. Whatever powers
the host is off-board and out of scope, with one caveat folded into Open
question 1: if the measured source is a battery pack above 30 V, the answer
changes this section.
\item \textbf{Motors: DOES NOT APPLY as specified}, with the same caveat. Measuring a
motor drive's current-sense or bus-divider output would put inductive
transients and a large common-mode component on the input of a board with no
protection. Covered by Open question 1.
\item \textbf{RF transmit: DOES NOT APPLY.} No radio function of any kind.
\item \textbf{High current (\textgreater{}3 A): DOES NOT APPLY.} The entire board draws an estimated
10 mA and carries no power path. The threshold is not approached.
\item \textbf{\textgreater{}30 V: DOES NOT APPLY at the stated numbers} - 5 V is the maximum on any
net, 3.3 V is the only rail, so \texttt{check\_creepage}'s derived spacing check is a
clean no-op. This too is contingent on Open question 1: a source that puts a
higher potential on the terminal, whether by fault or by design, engages it
for real.
\item \textbf{No input protection - noted, not flagged.} The scope tier excludes
protection of every kind, so an over-range input, a reversed input, or a
static discharge into the terminal can destroy the converter. On a bench
article with a low-voltage source that is an accepted consequence of the
declared scope, not a hazard to a person, and not a finding for a reviewer to
raise. It is recorded so the consequence is visible rather than implicit, and
it is why Open question 5 asks the owner to accept it knowingly.
\item \textbf{No isolation - noted, not flagged.} Input ground and host ground are the
same node. Any ground potential difference between the measured source and
the host appears directly across the input and is measured as signal.
Isolation is not in the scope tier and is not proposed; the owner should know
the input must share the host's ground reference.
\end{itemize}
\subsubsection*{9. Open questions}
Nine questions, all closed-form, each with a default. Question 1 is safety-relevant and has no safe default for the hazardous case - it must be answered, not defaulted past. The rest can be taken as the stated default if the owner has no preference.

\paragraph*{ANSWERED by the owner 2026-08-16 (P0 batch) - these are now requirements}
\begin{quote}\small\ttfamily\raggedright
\textbar{} Q \textbar{} Answer \textbar{} Effect \textbar{}\\
\textbar{}---\textbar{}---\textbar{}---\textbar{}\\
\textbar{} 1 \textbar{} Low-voltage bench source (owner ruled "simplest option", which is this) \textbar{} See the validity envelope below. RE-CONFIRMED AT H1. \textbar{}\\
\textbar{} 2 \textbar{} {*}{*}0.1 \% class{*}{*} \textbar{} +/-5 mV of a 5 V reading at 25 C; +/-12 mV across 0-50 C, at the terminal, uncalibrated \textbar{}\\
\textbar{} 3 \textbar{} {*}{*}SPI{*}{*} \textbar{} 10 kSa/s is a maximum, not gap-free streaming (sub-part at default) \textbar{}\\
\textbar{} 4 \textbar{} Single-ended to ground (default) \textbar{} One signal wire + one ground wire at J1 \textbar{}\\
\textbar{} 5 \textbar{} No protection ACCEPTED (default) \textbar{} Input stays strictly inside 0-5 V; no negative excursion to resolve \textbar{}\\
\textbar{} 6 \textbar{} Header per default \textbar{} Single-row 0.1 in male, one edge: 3V3, GND, then SPI, +1 if the part needs DRDY/CONVST \textbar{}\\
\textbar{} 7 \textbar{} Host rail per default \textbar{} \textgreater{}=50 mA available, 3.3 V +/-5 \%, tens of mV noise, no sub-3.0 V operation \textbar{}\\
\textbar{} 8 \textbar{} {*}{*}3.3 V only{*}{*} \textbar{} The 0-5 V span MUST be scaled into the 3.3 V domain; the attenuator is inside the error budget \textbar{}\\
\textbar{} 9 \textbar{} Source \textless{}=1 kohm, board \textgreater{}=100 kohm (default) - {*}{*}NARROWED AT P2, see below{*}{*} \textbar{} The board must not load the source appreciably \textbar{}
\end{quote}
\textbf{Q9 REVISED at P2 (orchestrator ruling, owner-delegated):} the contract is now \textbf{source \textless{}= 200 ohm} (narrowed 5x from the answered 1 kohm) and {*}{*}board presents 1.00 Mohm{*}{*} (strengthened 10x from the answered 100 kohm). Both numbers go on the silkscreen. Reason: as answered the two specs are jointly impossible at 0.1 \% - loading error is Rs/Rtot, so 1 kohm into 100 kohm is 1 \% = 50 mV against a 5 mV budget. Holding Rs \textless{}= 1 kohm instead is possible (a 5 Mohm 0.1 \% divider closes RSS at 4.27 mV) but publishes an 18.7 mV worst case against a 12 mV target and makes the board's dominant uncertainty an ESTIMATED surface-leakage term on a megohm node rather than a vendor-guaranteed maximum. This narrowing restricts what may be connected and is therefore a visible H1 item.

\textbf{Q1 validity envelope - the design is valid ONLY for a source of this kind.} Nothing anywhere in the source circuit above 30 V; no mains connection; no motor drive or other inductive load; no battery pack above 30 V; already referenced to the same ground as the host; lead about a metre or less. If the real source ever falls outside that envelope, section 8 is wrong and the board must not be connected to it: it has no isolation and no protection, and its input ground is the host's ground. The owner delegated this to the default rather than describing an actual source, so it is re-presented at H1 for explicit confirmation before any money or copper is committed.

{*}{*}Consequence of 2 + 8 together (the design problem this board actually teaches).{*}{*} A 0.1 \% class result must be delivered through an attenuator that brings 0-5 V into a 3.3 V domain, so the attenuator's ratio error and its tempco are first-class terms in the error budget alongside the reference. That combination - not the converter's bit count - is what P1/P2 have to solve.

1. {*}{*}What actually produces the 0-5 V signal you want to measure, and can any part of that circuit reach a dangerous or elevated voltage?{*}{*} This is the safety question. Say what it is in plain terms - a sensor, a potentiometer, the output of another bench instrument, a divider across a battery or a power supply, a current-shunt amplifier on a motor drive - and roughly how high the voltage goes ANYWHERE in that circuit, not just at the two wires coming to this board. This board has no isolation and no protection, and it shares the host's ground, so its terminal is electrically part of whatever it is connected to. DEFAULT (usable only if it is true): a low-voltage bench source - nothing in it above 30 V, no mains connection, no motor or other inductive load, no battery pack above 30 V - already referenced to the same ground as the host, connected by a lead of about a metre or less. If mains, a motor drive, or anything above 30 V is involved anywhere in that circuit, say so plainly. It does not necessarily stop the board, but it changes the requirements and this section has to be rewritten before the design proceeds. 2. \textbf{How accurate does "DC accuracy matters" actually have to be?} The brief says 12 bits or better, but bits are resolution, not accuracy - a 16-bit converter with a sloppy reference is less accurate than a 12-bit one with a good one, and this number is the single biggest driver of the parts on this board. Answer as an error at the terminal, with no user calibration assumed. Pick one: (a) \textbf{0.1 \% class (DEFAULT)} - total error within about +/-5 mV of a 5 V reading at 25 C, and within about +/-12 mV across the full 0-50 C range. An honest, ordinary instrument-grade number. (b) 1 \% class - about +/-50 mV. Cheaper, allows a converter's internal reference, and makes the board mostly about the conversion rather than the precision. (c) 0.01 \% class - about +/-0.5 mV. This is a precision-metrology requirement: it forces a high-grade external reference, 16 bits or more, and low-drift resistors, and it will dominate the parts cost. 3. \textbf{Which digital interface should the host speak: SPI or I2C?} DEFAULT: \textbf{SPI} - it is what precision converters in this class almost always offer, it costs no throughput at 10 kSa/s, and it needs no address management. Choose I2C only if the host has no free SPI port or already runs an I2C bus you want this board to join; at 10 kSa/s continuous, I2C would need to run at 400 kHz or faster. Related, answer only if it is not the default: is gap-free streaming at the full 10 kSa/s required, or is 10 kSa/s a maximum the host will rarely sustain? DEFAULT: a maximum, not an obligation. 4. {*}{*}Is the input a single voltage measured against ground, or a differential / floating measurement?{*}{*} DEFAULT: \textbf{single-ended} - one signal wire and one ground wire, the signal measured with respect to the same ground the host supplies. Answer "differential" only if the two wires from your source do not share the host's ground, or if you need to reject a voltage common to both. 5. {*}{*}What happens if the input goes outside 0 to 5 V - and do you accept that the board has no protection against it?{*}{*} There is no clamp, no series resistor beyond what the converter needs, and no reverse protection on this board, because the study scope excludes protection. A 12 V bench supply clipped onto the terminal by mistake, or the wires swapped, will most likely destroy the converter. DEFAULT: \textbf{accepted.} The board is rated 0 to 5 V, survives only what the chosen converter inherently survives, and nothing is added to protect it. Also DEFAULT: the signal stays strictly within 0 to 5 V in normal use - it never sits slightly below zero (as some sensors do at their zero point), so the design does not have to resolve negative inputs. Tell us if either default is wrong; a small negative excursion in particular changes the converter choice more than it sounds like it should. 6. \textbf{What should the host header look like physically?} Pin count, pin order and orientation were not stated, and the host presumably already exists. DEFAULT: a single-row 0.1 inch male pin header on one board edge, pins in the order 3.3 V, GND, then the digital signals for whichever interface answer 3 selects (4 signal pins for SPI, 2 for I2C), so 6 pins for SPI or 4 for I2C - plus one more pin if the chosen converter needs a data-ready or convert-start line. If your host expects a particular pinout or a particular connector, give it and it will be matched exactly. 7. {*}{*}How much current can the host's 3.3 V rail supply to this board, and how clean is it?{*}{*} The board's own draw is small (an estimated 10 mA), but if the design ends up measuring against that rail, its noise and tolerance become part of the measurement error. DEFAULT: at least 50 mA available, an ordinarily regulated 3.3 V +/-5 \% rail with tens of millivolts of noise on it, and no requirement for the board to work below 3.0 V. 8. {*}{*}Does the host header also have 5 V available, or is 3.3 V genuinely all there is?{*}{*} The brief says 3.3 V, and that is being taken at face value - but the input span is 0-5 V, so the answer changes the shape of the design more than any other question here. DEFAULT: \textbf{3.3 V only}, exactly as the brief states. The board will be designed to measure a 0-5 V input from a 3.3 V-only supply. 9. {*}{*}What is driving the terminal, electrically - can it supply a little current, or is it a high-impedance source that must not be loaded?{*}{*} A potentiometer wiper or a long-settling sensor output can be loaded down by the converter's own input, and the fix (a buffer) is a part on the board. DEFAULT: the source can drive at least a light load - source impedance of 1 kohm or less - and the board presents at least 100 kohm at the terminal. Tell us if your source is high-impedance (say, above 10 kohm) or if you need the board to be effectively invisible to it.

\section{Architecture}
\subsection*{Decisions of record (state.json)}
{\small
\begin{longtable}{lp{5.2cm}p{7.2cm}}
\toprule
\textbf{Phase} & \textbf{What} & \textbf{Why} \\
\midrule
\endhead
P0 & Build mode: mode block-basics; scope block-only; binding canonical & token 'learning block-basics:' in the brief; contract in reference/build-modes.md \\
P0 & Accuracy class = 0.1 percent at the terminal, uncalibrated: +/-5 mV of a 5 V reading at 25 C, +/-12 mV across 0-50 C & owner answer to P0 Q2; bits are resolution not accuracy, so this and not the 12-bit floor is the BOM driver \\
P0 & Host interface = SPI, 3.3 V CMOS, board is the peripheral; 10 kSa/s is a maximum, not gap-free streaming & owner answer to P0 Q3 \\
P0 & J2 supplies 3.3 V ONLY - no 5 V rail available; the 0-5 V span must be scaled into the 3.3 V domain on-board & owner answer to P0 Q8, brief taken at face value; makes the attenuator ratio error and tempco first-class terms in the 0.1 pct error budget \\
P0 & Safety envelope: source is a low-voltage bench source - nothing anywhere in it above 30 V, no mains, no motor drive, no \textgreater{}30 V pack, already sharing host ground, lead \textless{}=1 m. Design is valid ONLY inside this envelope & owner ruled 'simplest option' on P0 Q1, which resolves to the stated default; the board has no isolation and no protection and its input ground is the host ground, so the envelope is the safety argument. Re-confirmed at H1 before copper or money \\
P0 & Defaults taken for P0 Q4/Q5/Q6/Q7/Q9: single-ended to ground; no-protection accepted with no negative excursion; single-row 0.1in header 3V3-GND-SPI(+DRDY/CONVST if required); host rail \textgreater{}=50 mA 3.3 V +/-5 pct; source \textless{}=1 kohm and board presents \textgreater{}=100 kohm & owner reviewed the defaults in the P0 batch and did not override them \\
P0 & Owner delegated design decisions to the pipeline: 'make decisions yourself based on what will give the best learning for the basic system'. H1-H4 run as report-and-proceed (digest presented, verdict recorded as owner-delegated), not as blocking holds & explicit owner instruction 2026-08-16 mid-run. Scope of the delegation: engineering/design choices only. It does NOT relax any gate, the coverage checks, research, DFM, or the section-8 safety envelope - a change to the Q1 safety envelope re-blocks on the owner, and nothing spends money \\
P1 & P1 roster: 2x research-component-scout (adc; afe-support = vref/divider/buffer), research-reference-design (afe topology families T1-T4), research-power-architect (analog supply + reference architecture). research-interface-spec NOT spawned & SPI at 3.3 V CMOS into an unspecified host is a design decision, not a standards-bound interface, and the screw terminal and 0.1 in header bind no standard - there is no spec to research. The accuracy-critical unknowns are the attenuator, the reference and the SAR input drive, so the roster is weighted there \\
P1 & CLOSED: no 5 V second rail. The ADS8681 path (on-chip attenuator, native 0-5.12 V input, but needs \textgreater{}=4.5 V AVDD, so a boost converter) is rejected; the board stays 3.3 V-only with an external-reference SAR behind an attenuator & Three independent reasons. (1) MODE: scope tier block-only excludes a second rail the block does not need IN ORDER TO WORK - and the block demonstrably works at 3.3 V (MCP3201-class SAR, separate VREF pin), so a boost is a feature, not support. (2) OWNER: P0 Q8 answered 3.3 V only, explicitly. (3) LEARNING VALUE: the attenuator-plus-reference error budget IS the lesson of this board; a boost would delete that lesson and park a switching regulator beside a high-impedance precision node. Scout finding recorded: every \textgreater{}=16-bit-accuracy SAR in JLC stock needs \textgreater{}=4.5 V AVDD - nothing that accurate runs natively at 3.3 V \\
P1 & Converter lead candidate: MCP3201 class (1-channel 12-bit SAR, SOIC-8, SEPARATE VREF pin 0.25 V..VDD, SPI). MCP3202 is DISQUALIFIED despite better stock/price & MCP3202's 8-pin package ties VDD and VREF to one physical pin, so it cannot take a true sub-VDD external reference - the JLC catalog attribute claims 'External' and is misleading. Caught by reading DS21290F, not by parts\_search. MCP3204 is the second-source hedge on the same SAR core \\
P1 & AFE topology = T1/T3: precision attenuator -\textgreater{} low-Ibias RRIO buffer -\textgreater{} SAR with external precision reference. T2 (bare divider, no buffer) and T4 (ratiometric to the +3V3 host rail) are both DISQUALIFIED & T4 by arithmetic: the host rail is +/-5 pct (owner-answered Q7) = +/-165..250 mV of reference error, 20-50x the +/-5 mV budget, and it is NOT true ratiometric cancellation because the 0-5 V input is independent of the 3.3 V rail. T2 by the SAR acquisition equation t = (N+1){*}ln2{*}(Rs+Rss){*}Cs: Q9 demands the board present \textgreater{}=100 kohm at the terminal, which puts the divider Thevenin impedance in the tens of kohm, and at 12 bits that is marginal-to-failing against MCP3201's 1.5-clock sample window - before input leakage across that same impedance is counted. The buffer is what reconciles 'high-Z to the source' with 'low-Z and fast to the SAR'; it is support the operating point requires, not filtering the mode excludes \\
P1 & Extended-tier JLC parts are allowed for the ATTENUATOR (matched network) as well as for the converter and the reference; requirements.md section 6 is widened accordingly & P1 evidence makes the attenuator load-bearing for accuracy: two uncorrelated discrete resistors give ratio error \textasciitilde{}sqrt(2)x tolerance and ratio TCR error \textasciitilde{}sqrt(2)x TCR, and at 0.1 pct / 25 ppm per C that alone is \textasciitilde{}7 mV at 25 C and \textasciitilde{}4.4 mV more over 0-50 C - the whole budget, spent on two resistors. Under mode block-only COST is relaxable and ACCURACY is not, so the ruling is forced \\
P1 & Accuracy is claimed on an RSS basis: RSS \textasciitilde{}4.8 mV at 25 C meets the +/-5 mV target; the worst-case SUM is \textasciitilde{}10.6 mV (0.21 pct) and WILL be stated as such in the design doc and on the board's own claim. The board is not advertised as worst-case 0.1 pct & Owner delegated the ruling. RSS is the defensible convention for uncorrelated error terms (TI SNAA320B sec 4: real error sits closer to the RSS result) and every dominant term here - reference initial accuracy, divider ratio, converter INL, buffer offset - is genuinely independent. Worst-case closure would need 0.02 pct-class reference AND 0.02 pct-class divider, which is a different board. Honesty rule: both numbers get published; a claim of uncalibrated worst-case 0.1 pct would be false \\
P1 & HARD SELECTION CONSTRAINT for P3: the converter's specified analog input range must include 0 V (AVSS) in the configuration used, and the buffer's output must swing to it under the SAR's near-zero DC load & A PGA-input or headroom-limited converter needing 100-200 mV above AVSS cannot measure a 0 V input on a single 3.3 V rail; the only fixes are a negative rail or a level shift, both of which are second-rail mode boundaries. Selection avoids the boundary entirely. MCP3201's 0..VREF range satisfies it \\
P1 & Extended-tier allowance covers ALL FOUR accuracy-critical parts: converter, voltage reference, attenuator, buffer op-amp. requirements.md section 6 is widened to match & Scout evidence: no JLC Basic part in ANY of the four sub-functions meets the 0.1 pct class. The cheap defaults were quantified, not hand-waved - REF3320/REF3325/MCP1501/LM4040 each consume or exceed the whole 5 mV 25 C budget on initial tolerance alone, MCP6001 spends 90 pct of it on offset (4.5 mV) and TLV9001 1.6 mV. Under block-only, cost is relaxable and accuracy is not \\
P1 & Architecture is NON-ratiometric and confirmed as such: the reference serves the converter's VREF pin only, and the attenuator taps the raw 0-5 V input independently. Reference error and divider ratio error are INDEPENDENT terms and both stay in the budget & Scout OPEN 2 asked for confirmation because a ratiometric arrangement would cancel reference error out of the divider math and re-weight the rankings. It cannot apply here: the 0-5 V source is not excited by this board's reference or rail, so there is nothing to cancel against - the same reason T4 (ratiometric to +3V3) was rejected \\
P1 & P3 sourcing rule: for AD8605-pattern parts, ONLY the Analog-Devices-branded LCSC row is acceptable; clone-brand rows under an identical MPN pattern are refused & Scout pulled a clone datasheet (TECH PUBLIC) for an AD8605-pattern part specifying 5 mV MAX input offset - 30-300x looser than the genuine part, on the same pinout and part-number pattern. On this board 5 mV of offset IS the entire 25 C budget. JLC lists six clone brands alongside genuine ADI stock under one search term, so brand is a selection criterion here, not a preference \\
P2 & CONVERTER OVERRULED: MCP3201 is out, ADS8326IB class (16-bit) is in. MCP3201-B's own DS21290F specifies gain error +/-5 LSB MAX = 6.10 mV at the terminal - 1.22x the ENTIRE +/-5 mV budget - plus +/-3 LSB offset & P1's scout ranked on INL and never read the gain/offset rows. The mechanism is general and worth keeping: LSB-denominated DC error specs shrink 16x going from 12-bit to 16-bit, so a 12-bit part's '+/-5 LSB' is a far bigger voltage than a 16-bit part's. ADS8326IB specifies offset (+/-1 mV), gain (+/-16 LSB) and INL (+/-2.5 LSB) as MAXIMA - the only candidate found that does - and runs from 2.7-5.5 V, so it works on the 3.3 V rail \\
P2 & ATTENUATOR: five equal 200k 0.02 pct / 10 ppm resistors in ONE series string, tapped 3:2 -\textgreater{} K = 0.400, Rtot = 1.00 Mohm, VREF 2.048 V (ADR4520 class), full scale 5.120 V (2.4 pct headroom). Matched network REFUSED & Three independent reasons the network loses. (1) A quad's realisable ratios are \{0.2, 0.25, 1/3, 0.5, ...\}; K=0.5 with a 2.5 V reference puts full scale at EXACTLY 5.000 V, so the stated maximum input clips on a real unit - a functional refusal, not an error term. (2) Deep-stock LT5400 rows carry 1k/5k elements, capping Rtot at \textasciitilde{}12 kohm = 417 mV of loading error. (3) The 100 kohm quad is 11 pcs at 61 USD. The equal-string wins on arithmetic too: n equal elements AVERAGE their tolerances, sigma\_K/K = (1-K){*}sqrt(1/n\_t+1/n\_b){*}t = 0.548t here versus 0.849t for an unequal pair - the sqrt(2){*}tol convention P1 used overstates the pair by 2.6x because it drops the (1-K) factor. Live parts\_search evidence: 0.02 pct grade stops at 200 kohm and 0.01 pct/2 ppm stops at 10 kohm in stock, so a megohm divider is NECESSARILY a lower-grade divider unless it is a series string \\
P2 & ACCEPTED (orchestrator ruling on the architect's H1 question): the board contracts to Rs \textless{}= 200 ohm at the terminal, narrowing owner-answered Q9 from \textless{}= 1 kohm, and presents 1.00 Mohm, strengthening it from \textgreater{}= 100 kohm. Both numbers go on the SILKSCREEN & Config C (keep Rs \textless{}= 1 kohm, 5 Mohm divider, 0.1 pct resistors) also closes RSS at 4.27 mV, but it publishes an 18.7 mV worst case against a 12 mV target and, worse, moves the board's dominant uncertainty from datasheet MAXIMA to an estimated surface-leakage term on a megohm node. The chosen config keeps every dominant term a number a vendor guarantees. This NARROWS the owner's stated envelope, so it is a visible H1 item, not an implementation detail - and the contract belongs on the silk because input impedance is a specification of the instrument \\
P2 & Stackup JLC2313\_1.6, 2 layers, 1 oz, HASL, no impedance control. Flat root schematic sheet, NOT hierarchical & Single-sided assembly hands the entire B.Cu to one unbroken GND pour, which is exactly what the analog-return constraint D1 demands; a 4th layer would add nothing this board can use. Flat sheet per LEARNINGS 2026-07-22: a child-sheet label exports as /\textless{}sheet\textgreater{}/\textless{}LABEL\textgreater{} and silently unhooks every constraints.json entry - the P4 defect that cost lumina-par, and why bb-buck and sbuck-5v3a are flat. 18 parts, 12 nets, one A4 page \\
P2 & Research budget: KEEP all four declared blocks (resistive-attenuator, precision-buffer, sar-adc, series-voltage-reference) as separate coverage slots - the architect offered to fold precision-buffer into sar-adc to save a task; declined & Topology records seed the library more durably than part records, and 'op-amp driving a SAR input' is a reusable unit standing alone. 4 of 6 tasks at P2 leaves 2 for P3's three ICs. If P3 opens a third gap, the per\_run cap is RAISED by a recorded decision rather than a slot being silently truncated - the cap exists to make spend visible, not to cap knowledge on the one board built to produce it \\
P2 & The 5-element equal string's PRIMARY justification is restated: 0.02 pct-grade resistance stops at 200 kohm in stock and the design needs Rtot = 1.00 Mohm, so five elements is the only way to reach the value at that grade. The tolerance-AVERAGING advantage is a secondary, statistical bonus and is NOT claimed as sourced knowledge & The attenuator researcher looked for and could not find any vendor publishing the series-combining statistics sigma\_K/K = (1-K){*}sqrt(1/n\_t+1/n\_b){*}t - Vishay states only the qualitative principle - and correctly refused to put it in a record. The design does not depend on it: WORST-CASE ratio error is (1-K){*}2t = 0.024 pct for the 5-string and 0.024 pct for an unequal pair alike, so the string never loses on worst case, and the P2 budget's worst-case column already assumes no averaging. Averaging improves only the RSS column. Recording this so the design doc does not cite as vendor knowledge something that is our own arithmetic \\
P2 & NEW LAYOUT CONSTRAINT (P6/P7, first order): ground offset between the attenuator string's bottom node and the converter's negative reference is NOT attenuated. Signal is divided by K=0.4; a ground offset passes at unity, so 1 mV of offset = 2.5 mV input-referred = half the terminal error budget. To be folded into constraints.json before P5 & Found by the resistive-attenuator research (pending second-reader verification). It converts an ordinary-looking copper decision into the single largest layout-controlled error term on the board, and it is invisible in a schematic - exactly the class of defect this board exists to teach. Consolidated constraints.json update happens once, after all four research tasks close, so P5 board setup reads the final set \\
P2 & ARCHITECTURE: the converter's pseudo-differential -IN is wired as a dedicated REMOTE-SENSE line to the attenuator string's BOTTOM NODE - not to a convenient ground point. This is a schematic requirement at P4 and a routing constraint at P7 & It cancels the unattenuated ground-offset error exactly, which is otherwise the largest layout-controlled term on the board. Algebra: with the tap at V\_tap referred to the string bottom GND\_A, +IN referred to the converter ground GND\_C carries V\_tap + (GND\_A - GND\_C) + Vos, and -IN tied to GND\_A carries (GND\_A - GND\_C); the difference is V\_tap + Vos and every ground offset drops out. Two independent research tasks produced the halves of this - the attenuator task found that a divider passes ground offset at UNITY while dividing signal by K (so 1 mV = 2.5 mV input-referred at K=0.4), and the sar-adc task found that the ADS8326 specifies -IN as a remote-ground sense. Neither task saw the other's output \\
P2 & OPEN RISK for P3/P4, recorded not resolved: the ADS8326 datasheet calls for a 47 uF bypass at REF and warns 'make sure the op amp can drive the bypass capacitor without oscillation', while precision series references typically specify a PERMITTED CAPACITANCE RANGE, not a minimum. 47 uF may sit outside the ADR4520-class stability window & This is a direct conflict between two parts' datasheets and it decides whether the board needs a reference BUFFER or a series isolation resistor - a part count change, so it must be settled before P4 schematic, not discovered at bring-up. The series-voltage-reference research task was already briefed to get the permitted range and the 'buffer needed' criterion; if it comes back unresolved this becomes a P3 datasheet-extraction requirement \\
P2 & REFERENCE RISK CLOSED from the datasheet spec table: ADR4520 conditions its whole spec on VIN 3-15 V with DROPOUT 1 V MAX over -40..125 C at no load AND at 2 mA, so the guaranteed floor is 2.048 + 1.0 = 3.048 V. At the 3.135 V worst-case rail the part is IN SPEC with 87 mV of margin. The family headline 'low dropout 300 mV at 2 mA' is footnoted (VOUT \textgreater{}= 3 V) and does NOT apply to the 2.048 V member & This also retroactively kills P1's pick: ADR4525A is 2.5 V, VOUT \textless{} 3 V, so it carries the same 1 V dropout and needs 3.50 V - ABOVE the 3.135 V worst-case rail. It could not have worked, and the P1 scout's '300 mV dropout leaves 335 mV headroom' came from the footnoted headline. 2.048 V is therefore not a preference but the HIGHEST usable reference voltage from this family on a 3.3 V rail, which independently confirms the architect's choice. Note also that dropout is DEFINED (p29) as the differential at which VOUT has already degraded 0.1 pct = 2.05 mV = 41 pct of the terminal budget, so the 87 mV is margin to a degraded point, not to a cliff \\
P2 & OPEN RISK CLOSED - no reference buffer and no series isolation resistor needed: ADR4520 specifies a stable output capacitance WINDOW of 1-100 uF, and the ADS8326's required 47 uF REF bypass sits inside it. The reference drives the SAR's reference bypass directly & The conflict recorded earlier was that the ADS8326 datasheet warns 'make sure the op amp can drive the bypass capacitor without oscillation' while precision references specify a permitted RANGE rather than a minimum. Both halves are now sourced: the range is 1-100 uF and 47 uF is inside it. Part count unchanged, and the P4 schematic can proceed \\
P2 & ASSEMBLY CONSTRAINT: the voltage reference must be soldered in the LAST reflow pass. Solder-heat shift is +/-0.02 pct (ADI spec row; TI measured +0.03 pct LCCC / -0.04 pct VSSOP) - as large as the entire initial-accuracy spec - and is only removed by calibration, which this board does not do & Recorded now so P9/P10 cannot lose it. This board is single-sided SMT assembly, so there is exactly ONE reflow pass and the constraint is satisfied by construction - but it becomes live the moment anyone adds a second side or a rework step, and it is invisible in the schematic \\
P2 & BUFFER: OPA333 FAILS the settling requirement at the declared t\_acq of 9 us, on three independent numbers. Baseline moves to an OPA320-class fast RRIO CMOS part so P3 sourcing is unblocked; OPA333 stays a live alternate CONTINGENT on the t\_acq question below & TI's own chain gives a driver GBW floor of 4/(2{*}pi{*}R\_FLT{*}C\_FLT) with R\_FLT \textless{}= t\_acq/(12{*}C\_FLT), so the floor collapses to GBW \textgreater{}= 48/(2{*}pi{*}t\_acq) - INDEPENDENT of R and C, you cannot buy your way out with a bigger reservoir cap. At t\_acq = 9 us that is 849 kHz against OPA333's 350 kHz. Corroborated by OPA333's own settling curve (\textasciitilde{}40-50 us for a 4 V step) and by slew alone (0.16 V/us spends 12.8 us on a 2 V step). Cost of the swap: OPA320's Vos 150 uV max and 5 uV/C drift versus OPA333's 10 uV and 0.05 uV/C, i.e. \textasciitilde{}0.4 mV at 25 C and \textasciitilde{}0.7 mV over temperature at the terminal - affordable against a 3.23 mV RSS, and it moves RSS only to \textasciitilde{}3.30 mV \\
P3 & P3 VERIFICATION ITEM, decides the buffer: is the ADS8326's acquisition time HOST-CONTROLLED (the CS-high interval, so it stretches at low sample rates) or fixed by a clock-cycle count? At 10 kSa/s the sample period is 100 us, so if acquisition is host-controlled it can be 20 us+, the GBW floor drops below 350 kHz and OPA333 becomes viable and PREFERRED (30x lower drift). If it is clock-fixed at \textasciitilde{}9 us, OPA320 stands & The GBW floor is 48/(2{*}pi{*}t\_acq), so this single datasheet fact moves the requirement by a factor of \textasciitilde{}2.4 and changes which amplifier the board carries. It is a reading of the timing section, not a simulation result, so it belongs to datasheet extraction. The sim bench acquisition-settling then confirms whichever part the answer selects \\
P2 & OPA333's 240 kohm source impedance is NOT a problem if it is chosen: TI's zero-drift selection table gives OPA333 a 1000 kohm max recommended source impedance, and the divider's Thevenin is 240 kohm - inside it. The zero-drift family trades that ceiling against GBW in opposite directions (OPA189: 14 MHz but only 1 kohm) & Recorded so the OPA333-vs-OPA320 decision is made on SPEED alone, which is where the real conflict is, rather than being confused with a source-impedance objection that does not apply. Also flagged for the sim/firmware note: a chopper's clock tone can alias into a sampled system - OPA333 chops at \textasciitilde{}125 kHz which at 10 kSa/s folds to 5 kHz rather than to DC, but a sample rate that makes f\_chop an exact multiple would fold it to DC as a systematic offset \\
P2 & The remote-sense architecture is CONFIRMED against the source, not just inferred: TI SPRAD89 s4.2 states verbatim 'The common ground noise will not be attenuated by the resistor divider circuit', and s4.1 eq 8 independently multiplies a downstream offset by (R1+R2)/R1 = 1/K. Second reader re-read both pages and confirmed the arithmetic direction & This is the claim the board bet its -IN wiring and a P7 routing constraint on, so it was the primary refutation target handed to the fresh-context reader. It survived, twice over, from the source's own conclusions rather than from a plausible mechanism \\
P2 & The -IN remote-sense line is BOUNDED: ADS8326 specifies the -IN common-mode window as -0.3 V to +0.5 V relative to the device ground (datasheet p20, repeated p21, verified by second reader). The sense line to the attenuator's bottom node must stay inside it & This is the constraint that makes the remote-sense architecture safe rather than open-ended. Expected ground offsets on this board are millivolts, three orders inside the window, so the architecture is comfortable - but the bound must be recorded because it is what stops a future revision from running the sense line to a distant or noisy ground and calling it Kelvin sensing. P3 caveat: ADS8326 fig 39 (common-mode range vs VREF) is plotted at VDD 5 V with the VREF axis starting at 2.5 V, so this board's 3.3 V / 2.048 V point is off the plotted range - confirm the window at the real operating point during datasheet extraction \\
P2 & THIRD independent reason the OPA333 is rejected, and the cleanest one: the two constraints on the series resistor R\_FLT have an EMPTY intersection at 350 kHz. TI's settling ceiling (R \textless{}= t\_acq/12C, sprpeo2 p27) gives R \textless{}= 750 ohm at 1 nF (\textless{}= 1.6 kohm at 470 pF), while TI's stability optimum (sboa418 eq 2, with OPA333's R\_O = 2 kohm and f\_gbw = 350 kHz) asks for \textasciitilde{}1.2 kohm (\textasciitilde{}2.0 kohm at 470 pF). The stability optimum EXCEEDS the settling ceiling at both capacitor values - no legal R exists & Found by the fresh-context second reader while verifying both formulas on their own pages. It is stronger than the GBW-floor argument because it needs no extrapolation: two vendor formulas, both verbatim on cited pages, evaluated at this board's operating point, and they do not overlap. An amplifier can be REJECTED because the settling-derived ceiling and the stability-derived optimum on the same resistor have no common value - a general selection test worth carrying, and neither existing record names it. Requested as a new record \\
P2 & Flagged for the promotion ruling, not a defect: the attenuator ratio PRINCIPLE record says ratio error is set by relative specs 'not by their absolute tolerance or absolute TCR', while its TCR specialization says absolute TCR magnitude still moves the ratio. Both are Vishay's own words in DIFFERENT documents (AN 28194 vs TN 60108). At approval the parent's headline should be softened to match the specialization, which carries the nuance the parent elides & Surfaced by the second reader after the corrected record restored the qualifier 'but even in this case the ratio of resistance values is influenced by the magnitude of their TCR'. It is a genuine tension between two vendor documents, not a misreading, and it matters practically: a designer who takes the parent's headline alone concludes that buying a tracking spec ends the analysis, when the absolute TCR of the chosen network still has to be checked \\
P2 & PROMOTION IS DEFERRED to a separate owner pass AFTER the parallel board runs finish. This run will close with verified records reading  at coverage, which is the contract's expected state, and will NOT copy records into reference/knowledge/ during the run & research.py promote copies records and their quarantined sources INTO the shared library outside this workspace, and sibling board sessions (bb-ldo, bb-amp, bb-mcu) are running concurrently against the same git index and the same library directory. Concurrent promotion is how two runs clobber each other's library writes. Designing under  is explicitly allowed - the recipe says research narrows a gap to provisional and the owner closes it - so nothing in this board is blocked by deferring \\
P2 & CORRECTION to the preceding decision, whose text lost one word to shell quoting: the deferred-promotion decision should read "will close with verified records reading PROVISIONAL at coverage". No other word was affected and the ruling is unchanged & Backticks in the message were interpreted by the shell before state.py saw them. Recording the correction rather than leaving a sentence that reads as if the coverage state were blank, since the whole point of that decision is which state the run closes in \\
P2 & CORRECTION to the earlier reference decision. ADR4525 (2.5 V) is NOT disqualified on headroom: its dropout row is 500 mV MAX at both no load and 2 mA, so its floor is 3.0 V (set by the tables VIN condition) and it has 135 mV of margin on a 3.135 V rail - MORE than the ADR4520s 87 mV. The claim that it "would need 3.50 V" was wrong: it applied the 2.048 V members 1 V dropout to the 2.5 V member & Caught by the fresh-context second reader, which read the sibling members rows and listed the per-member floors. Dropout in this family is per-member and NOT monotonic in output voltage - 1 V at 2.048 V, 500 mV at 2.5 V, 100/300 mV at 3.0 V - so the 2.048 V part has the WORST dropout of the family, which is the opposite of the usual intuition. The footnoted-headline finding stands (300 mV at 2 mA is qualified to VOUT \textgreater{}= 3 V and applies to neither member); only the per-member number was misapplied \\
P2 & ADR4520 at 2.048 V is CONFIRMED as the choice, on the corrected rationale: RANGE UTILISATION with a realisable equal-element ratio, not headroom. With K = 0.4 from a 5-element 3:2 string, VREF 2.048 V puts full scale at 5.12 V and a 5 V input at 97.7 pct of range. Reaching comparable headroom with a 2.5 V reference needs K \textasciitilde{} 0.4545 = 5/11 or 4/9, i.e. NINE TO ELEVEN equal elements, or K = 0.5 with two, which returns full scale to exactly 5.000 V and clips the stated maximum input & Equal-element strings can only realise ratios n\_b/(n\_t+n\_b), so the reference voltage and the element count are one choice. 2.048 V is the value that lands a usable ratio at a sensible part count. Second-order confirmation: at FIXED full scale a higher VREF refers the converters OFFSET to the terminal by a smaller 1/K (2.05 vs 2.5 for a 1 mV offset), which is a \textasciitilde{}0.45 mV improvement - real but not worth six more resistors, and INL at the terminal is unchanged because VREF/K = FS \\
P2 & OPA333 rejection is now CLOSED, not merely argued: the obvious escape route is shut by a third vendor rule. Shrinking C\_FLT does not rescue it, because the settling ceiling falls as 1/C while the stability optimum falls only as 1/sqrt(C) - they cross at \textasciitilde{}243 pF, which is BELOW the C\_FLT \textgreater{}= 20{*}C\_SH = 960 pF floor the same TI slide imposes for a 48 pF sampling cap. Quantified remedy: at 1 nF with R\_O = 2 kohm, f\_gbw must exceed \textasciitilde{}0.78 MHz before ANY legal series resistor exists & Three TI rules interlock - the settling ceiling, the stability optimum, and the 20x charge-bucket floor - so there is no capacitor value at which a 350 kHz amplifier works here. Independently corroborated by OPA333 p11 fig 15 showing \textasciitilde{}40 pct overshoot at 1 nF, i.e. the part really is under-damped at this load. Caveat carried in the record and not hidden: sboa418 eq 2 assumes a FLAT resistive open-loop output impedance across the band, while OPA333 publishes only a single-frequency R\_O of 2 kohm at 350 kHz and no Z\_O-vs-frequency curve. The quoted frequency is f\_gbw, which is the one that matters, but flatness is unverified \\
P3 & BUFFER SELECTION TEST to apply at P3, generalised from the above: for any candidate amplifier, compute BOTH the settling ceiling R \textless{}= t\_acq/(12{*}C\_FLT) AND the stability optimum from R\_O and f\_gbw, and reject the part if the interval is empty at every C\_FLT above the 20{*}C\_SH floor. This is a cheaper and more decisive test than settling simulation and it runs on datasheet numbers alone & It is the test that killed the OPA333 without a single simulation, and it will re-run identically against whatever P3 sources. Recording it as a rule rather than a result so the P3 part-sourcer applies it to every candidate instead of only re-checking the incumbent \\
P2 & Source-impedance knee located from TI plot D039, measured off the plot frame rather than described: ENOB is flat and the curves are indistinguishable out to \textasciitilde{}3.3 kohm (11.94 at 3.3 k, 11.91 at 330 ohm, 11.86 at 33 ohm - the 3.3 kohm curve is the HIGHEST), and the knee falls between 3.3 kohm and 10 kohm (11.00 at 10 k, 9.59 at 20 k). A record had claimed 3.3 kohm \textasciitilde{}11.8, inventing a degradation in the low kohm and inverting the curve order & It is the record whose whole job is pricing source impedance, so a wrong curve reading there is load-bearing. The corrected shape changes the engineering story: moderate source impedance costs nothing, and the penalty arrives abruptly rather than gradually. This board is unaffected either way - its divider Thevenin is 240 kohm, an order past the knee, so the buffer is required regardless - but the number would have been wrong in the library \\
P2 & CORRECTION: the "empty intersection" argument OVERREACHES and is withdrawn as a rejection test. sboa418 p9 names 45 deg as generally stable and explicitly permits a SMALLER series resistor than eq 2 when voltage error dominates - eq 2 gives an OPTIMUM, not a MAXIMUM - so a ceiling below the optimum does NOT mean no legal R exists. At 470 pF the ceiling is 82 pct of the optimum, inside the 50-60 deg band eq 2 targets. The record carrying this test is refuted on that point & Caught by a fresh reader that re-derived all four algebraic claims (ceilings 750 / 1595.7 ohm, optima 1207.8 / 1956.4 ohm, crossover 243.5 pF, minimum 778 kHz f\_gbw) and found the ARITHMETIC correct but the INFERENCE unsupported by the source. The records own source note already conceded eq 2 is an optimum not a maximum, and the rule field then contradicted it. My earlier characterisation of the OPA333 rejection as "closed, not merely argued" rested on this leg and was wrong \\
P2 & OPA333 REJECTION STANDS on two independent legs that survive, neither involving eq 2: (1) the settling ceiling R \textless{}= t\_acq/(12{*}C) combined with the driver requirement GBW \textgreater{} 4/(2{*}pi{*}R{*}C) collapses to GBW \textgreater{}= 48/(2{*}pi{*}t\_acq) = 849 kHz at t\_acq = 9 us, against OPA333 350 kHz; (2) OPA333 fig 14 measured at unity gain settles in 40.4 us to 0.01 pct and 50.5 us to 0.001 pct, against a 9 us window. Baseline stays OPA320-class & Both legs are sourced and verified independently of the withdrawn collision test - leg 1 is arithmetic on two separately verified TI rules, leg 2 is a direct datasheet measurement of the part itself. The conclusion is unchanged; only one of the three supporting arguments fell. Note this rejection is still CONDITIONAL on t\_acq = 9 us, which remains the open P3 question - if acquisition proves host-controlled and stretches past \textasciitilde{}22 us, leg 1 dissolves and leg 2 must be re-judged against the per-conversion event rather than a full-scale step \\
P2 & PROCESS FINDING for the review loop: a non-binding "blemish" note from one second reader caused a CORRECT record to be edited into a WRONG one. The earlier reader flagged fig 15 as reading \textasciitilde{}7 pct / \textasciitilde{}15 pct without refuting; the researcher applied it; the fresh reader then measured the curve geometry and found the ORIGINAL \textasciitilde{}10 pct / \textasciitilde{}20 pct was accurate (true values 9.8 pct at 100 pF, 24.6 pct at 300 pF - and the curve does not extend to 1 nF at all, spanning only 10.4-604.5 pF) & Flag-without-refute has a real cost: an advisory note carries the authority of the reviewer role without the evidentiary bar of a refutation, and a researcher will act on it. Either raise it to a refutation with the measurement behind it, or mark it explicitly as do-not-act-without-verifying. Worth carrying into the skill, not just this board \\
P2 & STRONGEST and UNCONDITIONAL argument for the buffer, independent of t\_acq and of any amplifier choice: a switched-capacitor SAR input presents an equivalent average resistance R\_eq = 1/(f\_sample {*} C\_sample). At 10 kSa/s with C\_SH = 48 pF that is \textasciitilde{}2.08 Mohm. The attenuators 240 kohm Thevenin driving \textasciitilde{}2.08 Mohm is a gain error of 240k/(240k+2080k) = \textasciitilde{}10 pct - about 500 mV at the terminal, a hundred times the budget. A bufferless divider is dead on DC loading alone, before settling is discussed & The two existing arguments are both conditional - the GBW floor depends on t\_acq = 9 us (still an open P3 question) and the 40 us settling figure is specific to the OPA333. This one depends on neither: it holds at any acquisition time and for any converter with a capacitive sampling input, because slowing the sample rate raises R\_eq rather than lowering it. Mechanism is TI SPRAD89 equivalent-average-resistance, surfaced by the resistive-attenuator research; the arithmetic here is ours. It means the buffer is a REQUIREMENT of the topology, and only WHICH buffer remains open \\
P2 & ENOB-vs-source-impedance is a function of SAMPLE RATE, not source impedance alone: on TI plot D039 at 10 kSPS all five curves including 20 kohm sit at 11.9-12.0, while at 100 kSPS the knee falls between 3.3 kohm and 10 kohm. The knee is recorded scoped to the 100 kSPS end with the plots conditions (14-bit converter, CFLT 680 pF, 2-100 kSPS sweep), not as an impedance constant & Recorded because it is the correct general statement and because it looks, at first glance, like it might rescue a bufferless design on THIS board - we run at 10 kSa/s, where the plot shows impedance barely matters. It does not rescue it: 240 kohm is 12x beyond the plots 20 kohm maximum, and the equivalent-average-resistance argument above kills it independently at any rate \\
P2 & The OPA333-vs-OPA320 choice is formally handed to the SIM LEG, and that is the sources own instruction: TI sboa418 p11 states "Spice simulation tools such as TINA-TI must be used to verify the resultant phase margin" and publishes NO closed form for the 45-degree isolation resistor. The collision test therefore SCREENS a candidate and triggers a phase-margin simulation; it cannot reject one on paper & This closes the loop honestly. The screening test still earns its place - it flags, from datasheet numbers alone, that this amplifier must run an under-damped isolation resistor - but the vendor is explicit that the resulting margin is only knowable by simulation. The architect already named buffer-stability (PM \textgreater{}= 45 deg, overshoot \textless{}= 5 pct) and acquisition-settling (15.6 uV inside 9 us) as sim benches, so the machinery exists and the P8 sim gate is where this gets settled. The OPA320-class baseline stands meanwhile on the two legs that do not need simulation: the 849 kHz GBW floor at t\_acq = 9 us, and OPA333 fig 14s measured 40.4 us settling to 0.01 pct at unity gain \\
P2 & CLOSED sar-adc at 10/10 verified WITHOUT spending a further round on the one improvement the reader recommended: adding SBAA256B p.3 to the settling records sources. That page carries a TABLE giving 10.99 ENOB at RTH 10k / CFLT 680 pF / 100 ksps - a non-graphical, independent corroboration of the exact curve endpoint the record reads off a plot - plus the plots unstated conditions (100 Hz analog input, no driver amplifier). Logged as a promotion-time improvement instead & The record is already verified and its numbers were confirmed twice by two different measurement methods, so the citation adds robustness rather than correctness, and editing a verified record would void its signed attestation and cost a full re-open plus re-read. Proportionality: three correction rounds have already run on this task and the marginal finding no longer changes a rule. Whoever promotes these records into the library should add p.3 then, when the record is lawfully open anyway \\
P2 & DESIGNING UNDER A COVERAGE GAP, recorded not silent: slot block:B1 resistive-attenuator remains `gap` on the single class `constraints-emission`; all five of its other classes and all classes of B2/B3/B4 read `provisional`. No further research task is spent on it. The boards own constraints.json will be populated for this block from the six VERIFIED attenuator records directly & The attenuator researcher declared this class an honest gap rather than inventing a record for it, and the reader confirmed the checklist is honest about it. Proportionality: 4 tasks x 3 correction rounds have produced 32 verified records and 4 checklists, and 2 research tasks remain for P3s three ICs, where vendor layout sections are the higher-value target. The gap costs the LIBRARY a record that would tell future boards which machine-checkable constraints a precision attenuator implies; it costs THIS board nothing, because the underlying mechanisms - guard ring on the high-Z tap, thermal gradient across the string, the bottom-node sense - are all carried by verified records and go into constraints.json anyway. Flagged as the first candidate if the owner reopens research on this topology \\
P2 & REVERSED, on the datasheets own evidence: the converter DOES get a 5-10 ohm series resistor at its rail entry, with 0.1 uF + 10 uF at the pin, exactly as the ADS8326s own recommended application circuit shows (figs 44/45, p27). My earlier no-ferrite-no-series-R ruling is superseded for this element & The scope tier excludes "filtering the datasheet does not require" - it does NOT exclude what the datasheet DOES require. The P1 power architect wrote precisely this escape hatch ("it goes in if the chosen converters own datasheet shows it") and the condition is now met by a verified, page-cited record. This is the mode working as designed rather than being bent: a reflex ferrite gets refused, a vendor-specified element does not. Cost is one 0603 and, at this boards \textasciitilde{}60 uA typical draw at a 10 kHz data rate, \textasciitilde{}0.6 mV of drop. Separately, the AVDD/DVDD framing that motivated the original ruling is MOOT: the MSOP-8 ADS8326 has a single VDD pin, so there was never a split to bridge - the figs-44/45 element is rail-entry isolation, a different thing \\
P2 & ACCEPTED as an explicitly-labelled inference, not a page-cited rule: the reference keepout from mounting holes and the board edge. What IS page-cited is reflow package die-stress (+/-0.02 pct, as large as the whole B-grade initial accuracy) plus the vendors warning that its magnitude depends on board size, thickness and material. The step from that to "keep the part away from flex concentrators" is ours and is marked as such in the constraints trace table & The constraint costs nothing on a board this sparse and protects the single most accuracy-critical part, so it stays. But it must not masquerade as vendor knowledge in a library the whole run exists to seed - an inference labelled as an inference is honest; the same sentence unlabelled would be exactly the failure mode four correction rounds were spent removing \\
P4 & P4 SCHEMATIC HARD REQUIREMENT, flagged because the sheet plan almost lost it: U1 -IN is a DEDICATED SENSE RUN to the attenuator strings bottom node (R5s pad), NOT a GND symbol placed at the converter. sheets.md said "U1 -IN = GND", which is electrically true and operationally fatal & A P4 agent reading "= GND" wires it to the nearest ground and silently deletes the boards largest error-cancellation mechanism - the thing that turns a 2.5 mV input-referred ground offset into zero. The schematic looks identical either way and no gate would catch it: ERC sees one net either way, DRC sees copper either way, and the board would simply read wrong by half its error budget. This is the exact class of defect that only a stated requirement prevents \\
\bottomrule
\end{longtable}
}
\subsection*{Sheet plan}
\subsection*{sheets.md - bb-adc sheet plan, refdes ranges, CANONICAL net names}
\subsubsection*{1. Sheet plan: ONE FLAT ROOT SHEET}
\begin{quote}\small\ttfamily\raggedright
\textbar{} Sheet \textbar{} File \textbar{} Blocks it carries \textbar{} Interface nets at its boundary \textbar{} Refdes ranges \textbar{} `pwr\_base` \textbar{}\\
\textbar{}---\textbar{}---\textbar{}---\textbar{}---\textbar{}---\textbar{}---\textbar{}\\
\textbar{} root (only) \textbar{} `kicad/bb-adc.kicad\_sch` \textbar{} B1 attenuator, B2 buffer, B3 converter, B4 reference, J1 input, J2 host header, H1-H4 mounting \textbar{} none - the board's only interfaces are the two physical connectors (`/AIN\_RAW` + `GND` at J1; `+3V3`, `GND`, `/CS`, `/SCLK`, `/DOUT` at J2) \textbar{} `U1-U9`, `R1-R19`, `C1-C19`, `J1-J9`, `H1-H9` \textbar{} {*}{*}1{*}{*} (`\#PWR01..`, `\#FLG01..`) \textbar{}
\end{quote}
{*}{*}A hierarchical split is deliberately REJECTED, and this is a considered deviation from the architect role's default instruction.{*}{*} The reason is recorded, machine-verified and has already cost two boards on this project: a wiring label inside a CHILD sheet exports to the netlist as \texttt{/\textless{}sheet\textgreater{}/\textless{}LABEL\textgreater{}}, not \texttt{/\textless{}LABEL\textgreater{}} (LEARNINGS 2026-07-22, kicad-sch-api 0.5.6 hierarchy semantics: child-internal nets become \texttt{/\textless{}sheet\textgreater{}/NAME}; only nets merged to the root through a sheet pin keep the root-side name). Every \texttt{constraints.json} entry that spells the net \texttt{/\textless{}LABEL\textgreater{}} then silently no-ops - \texttt{check\_current} finds no copper on a net that does not exist and reports nothing. That is the P4 amendment class that cost \texttt{lumina-par}, and the reason \texttt{bb-buck} and \texttt{sbuck-5v3a} are also flat.

This board has 20 electrical parts and 14 nets (s2's table is the count; the earlier "12" was already one short of its own table). It fits one A4 page with room to draw the guard ring symbolically. The cost of hierarchy here is real and its benefit is zero.

Refdes ranges are therefore trivially unique. They are stated anyway because the rule is "unique ACROSS sheets": {*}{*}if a later amendment ever splits this sheet, the natural cut is \texttt{afe} (J1, R1-R6, U3, C6, C7) and \texttt{digital} (U1, U2, R7, J2, C1-C5, C8), which take \texttt{pwr\_base = 40} and \texttt{80} and keep the prefix ranges above{*}{*} - and every net name in s2 must then be re-checked against the new export.

\subsubsection*{2. CANONICAL NET NAMES - these are now binding}
\texttt{research/power.json} proposed \texttt{+3V3 / VREF / GND} and left the signal names open. P2 fixes all of them here. These are the names the schematic MUST produce and the names \texttt{constraints.json} uses.

\begin{quote}\small\ttfamily\raggedright
\textbar{} Canonical net \textbar{} KiCad mechanism \textbar{} Spans \textbar{} Declared in constraints.json as \textbar{}\\
\textbar{}---\textbar{}---\textbar{}---\textbar{}---\textbar{}\\
\textbar{} {*}{*}`+3V3`{*}{*} \textbar{} power symbol (VALUE = net name) -\textgreater{} exports BARE \textbar{} J2.1 -\textgreater{} C1 -\textgreater{} {*}{*}R7{*}{*} , U2 IN, U3 V+ \textbar{} `power` 11 mA \textbar{}\\
\textbar{} {*}{*}`VDD\_ADC`{*}{*} \textbar{} power symbol with {*}{*}Value = `VDD\_ADC`{*}{*} -\textgreater{} exports BARE \textbar{} {*}{*}R7 -\textgreater{} C2 -\textgreater{} C8 -\textgreater{} U1 VDD.{*}{*} The converter's supply pin ONLY \textbar{} `power` 3 mA \textbar{}\\
\textbar{} {*}{*}`GND`{*}{*} \textbar{} power symbol -\textgreater{} exports BARE \textbar{} everything; B.Cu pour; J2.2 and J2.6 \textbar{} `power` 11 mA `plane\_fed` `pdn:false`, `planes` B.Cu \textbar{}\\
\textbar{} {*}{*}`VREF`{*}{*} \textbar{} power symbol with {*}{*}Value = `VREF`{*}{*} -\textgreater{} exports BARE \textbar{} U2 OUT -\textgreater{} C5 -\textgreater{} C3 -\textgreater{} U1 VREF pin \textbar{} `power` 1 mA \textbar{}\\
\textbar{} {*}{*}`/AIN\_RAW`{*}{*} \textbar{} root-sheet LOCAL LABEL -\textgreater{} ONE leading slash \textbar{} J1.1 -\textgreater{} R1 \textbar{} `voltages` 5 V \textbar{}\\
\textbar{} {*}{*}`/ATT\_A`{*}{*} \textbar{} root local label \textbar{} R1 -\textgreater{} R2 \textbar{} not declared (see s4) \textbar{}\\
\textbar{} {*}{*}`/ATT\_B`{*}{*} \textbar{} root local label \textbar{} R2 -\textgreater{} R3 \textbar{} not declared \textbar{}\\
\textbar{} {*}{*}`/AIN\_DIV`{*}{*} \textbar{} root local label \textbar{} R3 -\textgreater{} R4 -\textgreater{} U3 `+IN`. {*}{*}THE high-Z tap, 240 kohm Thevenin{*}{*} \textbar{} not declared - it is a placement/routing constraint, not a width rule (s4) \textbar{}\\
\textbar{} {*}{*}`/ATT\_C`{*}{*} \textbar{} root local label \textbar{} R4 -\textgreater{} R5 \textbar{} not declared \textbar{}\\
\textbar{} {*}{*}`/AIN\_BUF`{*}{*} \textbar{} root local label \textbar{} U3 OUT -\textgreater{} U3 `-IN` (follower) -\textgreater{} R6 -\textgreater{} {*}{*}and the GUARD RING copper{*}{*} \textbar{} not declared \textbar{}\\
\textbar{} {*}{*}`/AIN\_ADC`{*}{*} \textbar{} root local label \textbar{} R6 -\textgreater{} C7 -\textgreater{} U1 `+IN` \textbar{} not declared \textbar{}\\
\textbar{} {*}{*}`/CS`{*}{*} \textbar{} root local label \textbar{} J2.3 -\textgreater{} U1 CS (active low; the bar is documentation, not part of the net name) \textbar{} `high\_speed` ref GND \textbar{}\\
\textbar{} {*}{*}`/SCLK`{*}{*} \textbar{} root local label \textbar{} J2.4 -\textgreater{} U1 DCLOCK \textbar{} `high\_speed` ref GND \textbar{}\\
\textbar{} {*}{*}`/DOUT`{*}{*} \textbar{} root local label \textbar{} U1 DOUT -\textgreater{} J2.5 \textbar{} `high\_speed` ref GND \textbar{}
\end{quote}
Why \texttt{VREF} and not \texttt{+2V048}: it is a rail feeding a power pin and decoupled at both ends, so it is a power symbol, and a power symbol's exported name is bare. \texttt{VBUS} sets the precedent for a bare non-numeric power net. \texttt{VDD\_ADC} follows the same rule for the same reason - a rail feeding a power pin, bypassed at that pin - and it exists because {*}{*}R7 splits the converter's supply off \texttt{+3V3}\textbf{, and KiCad cannot put one net on both pins of a resistor. }A power symbol WINS over a coincident local label{*}{*} (LEARNINGS 2026-07-28), so P4 must not mix the two mechanisms on one node, and \texttt{component.value = x} is the only way to set a power symbol's Value (\texttt{set\_property("Value", ...)} is a silent no-op on kicad-sch-api 0.5.6) - which now applies to TWO symbols, \texttt{VREF} and \texttt{VDD\_ADC}.

\textbf{Enforcement, not hope:} at P4 run \texttt{netlist\_audit.py --net kicad/bb-adc.net --constraints architecture/constraints.json}. A net spelled differently in the schematic surfaces as \texttt{missing\_net} (error).

\subsubsection*{3. Contents of the root sheet}
\begin{quote}\small\ttfamily\raggedright
\textbar{} Refdes \textbar{} Part class \textbar{} Net attachments \textbar{}\\
\textbar{}---\textbar{}---\textbar{}---\textbar{}\\
\textbar{} `U1` \textbar{} 16-bit SAR, SPI peripheral, VDD 2.7-3.6 V, external VREF pin 0.1 V..VDD, input range includes 0 V, offset AND gain specified as MAXIMA ({*}{*}ADS8326IB{*}{*} class) \textbar{} VDD={*}{*}`VDD\_ADC`{*}{*} (behind R7), GND=`GND`, VREF=`VREF`, `+IN`=`/AIN\_ADC`, {*}{*}`-IN`=`GND` AT R5's BOTTOM PAD - read the note below before wiring it{*}{*}, CS=`/CS`, DCLOCK=`/SCLK`, DOUT=`/DOUT` \textbar{}\\
\textbar{} `U2` \textbar{} 2.048 V series voltage reference, \textless{}= 0.02 \% initial max, \textless{}= 2 ppm/degC max, Vin\_min \textless{}= 3.035 V ({*}{*}ADR4520B{*}{*} class) \textbar{} IN=`+3V3`, GND=`GND`, OUT=`VREF` \textbar{}\\
\textbar{} `U3` \textbar{} unity-gain follower, CMOS input, RRIO, Vos \textless{}= 100 uV max, Ib \textless{}= 400 pA over temp ({*}{*}OPA333{*}{*} class; {*}{*}OPA320{*}{*} the same-footprint alternate) \textbar{} V+=`+3V3`, V-=`GND`, `+IN`=`/AIN\_DIV`, `-IN`=`/AIN\_BUF`, OUT=`/AIN\_BUF` \textbar{}\\
\textbar{} `R1` `R2` `R3` \textbar{} attenuator TOP arm, 3 x 200 kohm 0.02 \% / 10 ppm, ONE part number with `R4`/`R5` \textbar{} `/AIN\_RAW`-`/ATT\_A`-`/ATT\_B`-`/AIN\_DIV` \textbar{}\\
\textbar{} `R4` `R5` \textbar{} attenuator BOTTOM arm, 2 x the same 200 kohm part \textbar{} `/AIN\_DIV`-`/ATT\_C`-`GND` \textbar{}\\
\textbar{} `R6` \textbar{} ADC input isolation, 20-100 ohm (sim benches 2/3 set it) \textbar{} `/AIN\_BUF` / `/AIN\_ADC` \textbar{}\\
\textbar{} `R7` \textbar{} {*}{*}converter rail-entry isolation, 5-10 ohm (10 ohm class), 0402/0603 thick film{*}{*} - per the ADS8326's own recommended application circuit \textbar{} `+3V3` / `VDD\_ADC` \textbar{}\\
\textbar{} `C1` \textbar{} 10 uF X7R bulk, {*}{*}at the J2 entry, not at the converter{*}{*} \textbar{} `+3V3` / `GND` \textbar{}\\
\textbar{} `C2` \textbar{} 100 nF X7R, converter VDD - {*}{*}the SMALLER of the pair, so it sits closest to the pin{*}{*} \textbar{} `VDD\_ADC` / `GND` \textbar{}\\
\textbar{} `C3` \textbar{} converter VREF bypass - {*}{*}47 uF class low-ESR{*}{*} per the ADS8326's reference-input section, which is inside the ADR4520's 1-100 uF stable window (see `C5`) \textbar{} `VREF` / `GND` \textbar{}\\
\textbar{} `C4` \textbar{} 100 nF X7R, reference IN \textbar{} `+3V3` / `GND` \textbar{}\\
\textbar{} `C5` \textbar{} reference OUT cap - {*}{*}MANDATORY and \textgreater{}= 1 uF.{*}{*} NOT a DNP candidate \textbar{} `VREF` / `GND` \textbar{}\\
\textbar{} `C6` \textbar{} 100 nF X7R, buffer V+ \textbar{} `+3V3` / `GND` \textbar{}\\
\textbar{} `C7` \textbar{} 1-2.2 nF {*}{*}C0G/NP0{*}{*} at the converter input pin \textbar{} `/AIN\_ADC` / `GND` \textbar{}\\
\textbar{} `C8` \textbar{} {*}{*}10 uF X7R at the converter VDD pin{*}{*}, beside `C2` and behind `R7` \textbar{} `VDD\_ADC` / `GND` \textbar{}\\
\textbar{} `J1` \textbar{} 2-pole 5.08 mm THT screw terminal, silk-marked SIG and GND \textbar{} 1=`/AIN\_RAW`, 2=`GND` \textbar{}\\
\textbar{} `J2` \textbar{} 6-pin 2.54 mm single-row male THT header \textbar{} 1=`+3V3`, 2=`GND`, 3=`/CS`, 4=`/SCLK`, 5=`/DOUT`, 6=`GND` \textbar{}\\
\textbar{} `H1`-`H4` \textbar{} M3 clearance holes, 3.2 mm \textbar{} mechanical (no net) \textbar{}
\end{quote}
\paragraph*{\texttt{U1 -IN} IS A DEDICATED SENSE RUN, NOT A GROUND SYMBOL}
Electrically \texttt{-IN} is on \texttt{GND}, so a netlist cannot tell the difference - which is exactly why this note exists. {*}{*}Wire \texttt{U1 -IN} to the point where \texttt{R5}'s bottom pad meets \texttt{GND}{*}{*} (the attenuator string's bottom node), as its own connection. Do NOT drop a \texttt{GND} power symbol next to the converter and call it done: that deletes the board's largest error-cancellation mechanism silently and nothing downstream will report it.

The reason: {*}{*}a divider passes a ground offset at UNITY while dividing the signal by K = 0.400{*}{*}, so any offset between the string's bottom node and the converter's negative reference is referred to the input multiplied by 1/K = 2.5. \textbf{1 mV of copper offset = 2.5 mV at the terminal}, half the 25 degC budget. Sensing at the string bottom cancels it exactly: \texttt{+IN} carries \texttt{V\_tap + (GND\_A - GND\_C)}, \texttt{-IN} carries \texttt{(GND\_A - GND\_C)}, difference \texttt{V\_tap}. \textbf{Bounded:} the ADS8326 specifies \texttt{-IN} at \textbf{-0.3 V to +0.5 V} relative to device ground, so the sense point must be a node whose offset stays in the millivolts - which R5's pad is, and which a distant or shared-return ground is not. The bottom of the string is a \emph{reference tie}, not a ground connection: it must not share copper with a return carrying other current. Carried in \texttt{constraints.json} as the \texttt{R5}-\textgreater{}\texttt{U1} corridor; see \texttt{blocks.md} s8.1.

\paragraph*{\texttt{R7} + \texttt{C2} + \texttt{C8}: the converter's rail entry, and it is datasheet-required}
\texttt{R7} (5-10 ohm) goes between \texttt{+3V3} and \texttt{U1}'s VDD pin, {*}{*}upstream of both caps{*}{*}, with \texttt{C2} (0.1 uF, smaller, closest to the pin) and \texttt{C8} (10 uF) on the \texttt{VDD\_ADC} side. This is the ADS8326's own recommended application circuit, so it is what the datasheet REQUIRES, not "filtering the datasheet does not require" - the distinction the scope tier actually draws. A SAR has no supply rejection at the instant that matters, because the spikes that hurt land just before the comparator latches.

\textbf{It isolates `U1` alone.} \texttt{U2} and \texttt{U3} stay on \texttt{+3V3} upstream of \texttt{R7}, for the same reason the source figures put the reference and the op amps on the analog rail rather than behind the converter's RC. Cost: \textasciitilde{}0.6 mV of DC drop at a \textasciitilde{}60 uA typical draw, which costs the error budget nothing (the conversion is referenced to \texttt{VREF}, not to VDD) and leaves VDD inside 2.7-3.6 V on the worst -case rail. \textbf{No ferrite anywhere} - that one is still the reflex addition nothing here requires.

\paragraph*{\texttt{C5} is mandatory, and it is why there is no reference buffer on the BOM}
The ADR4520 specifies a stable output-capacitance {*}{*}WINDOW of 1 uF min to 100 uF max{*}{*}, two-ended, and the 2.048 V member is one of the two that need the full 1 uF (the \textgreater{}= 3.0 V members need only 0.1 uF). The output cap is a compensation element inside the loop, so {*}{*}both ends bind and \texttt{C5} cannot be DNP{*}{*}. That same window contains the ADS8326's 47 uF-class REF bypass (\texttt{C3}), which is what closed the open risk that this board might need a reference buffer or a series isolation resistor between the reference and \texttt{C3}: \textbf{neither is fitted, and the reference drives the converter's bypass directly.} If a current-limiting resistor were ever needed there it goes between the reference SOURCE and \texttt{C3} - never between \texttt{C3} and the VREF pin, where it would sit in the per-bit pulse path.

\textbf{J2 is 6 pins, and the sixth is GND, not a spare.} The chosen converter's SPI is read-only (CS, DCLOCK, DOUT - no MOSI), so the block needs 5 pins. The answered Q6 default reserved 6. Spending the sixth on a second GND puts a return reference at BOTH ends of the digital group, which is the return-path requirement D3/D4 already impose, rather than leaving a floating pin. It is support for the interface, not a feature.

\textbf{No series damping on `/SCLK` or `/CS`, no pull-up on `/CS`.} Both are "conditioning the datasheet does not require" and the scope tier excludes them; \texttt{/CS} is driven by the host at all times the board is powered, and the one hot-plug hazard that a pull-up would not fix is already recorded and accepted (answered Q5).

\textbf{The follower's feedback is a WIRE.} \texttt{U3 -IN} ties directly to \texttt{U3 OUT}, no resistor. A gain-setting network there would put two more tolerances and two more leakage nodes into the budget for no benefit - the attenuation is already done, precisely, upstream.

\subsubsection*{4. Nets deliberately left OUT of \texttt{constraints.json}}
\texttt{/AIN\_RAW} is declared only in \texttt{voltages} (5 V, which makes \texttt{check\_creepage} a clean, visible no-op). \texttt{/ATT\_A}, \texttt{/ATT\_B}, \texttt{/AIN\_DIV}, \texttt{/ATT\_C}, \texttt{/AIN\_BUF} and \texttt{/AIN\_ADC} are declared nowhere, on purpose:

\begin{itemize}
\item A \texttt{power} entry on any of them would give a high-impedance sense node a
minimum-width rule and pull it into \texttt{check\_pdn}'s decoupling inventory.
\texttt{/AIN\_DIV} in particular wants to be SHORT and THIN and nowhere near
anything, which is the opposite of what a width rule expresses.
\item A \texttt{high\_speed} entry would be a lie: they are DC nodes. Their sensitivity is
to LEAKAGE and to CAPACITIVE COUPLING FROM the SPI nets, which is why the
three SPI nets ARE declared \texttt{high\_speed} (that check verifies THEIR return
path is continuous, which is the mechanism that keeps their return current
out of the analog region) and why the protections for the analog nodes live
in \texttt{placement.corridors}, \texttt{placement.separation} and the prose rules D1-D5.
\end{itemize}
The rules that DO apply to them are layout rules, carried in \texttt{constraints.json} \texttt{placement} and in s5 below.

\subsubsection*{5. Placement groups this sheet implies (they become P6 groups)}
\begin{itemize}
\item \textbf{`attenuator`} (anchor \texttt{U3}; \texttt{R1}-\texttt{R5}, \texttt{R6}, \texttt{C6}, \texttt{C7}): the five string
elements in physical order R1-\textgreater{}R5 with the R3/R4 junction adjacent to \texttt{U3}'s
\texttt{+IN} pad. \texttt{/AIN\_DIV} is the SHORTEST high-impedance run on the board - a few
millimetres, not the most convenient route. {*}{*}The guard ring is part of this
group's geometry:{*}{*} \texttt{/AIN\_BUF} copper encircles the \texttt{/AIN\_DIV} trace and
\texttt{U3 +IN} on F.Cu with clearance on both sides, and solder-mask relief over
the ring. An annealer optimising a cost function cannot discover a guard
ring; this cluster is placed by explicit \texttt{place\_edit} and LOCKED before
\texttt{place\_anneal} runs.
\item \textbf{`reference`} (anchor \texttt{U2}; \texttt{C4}, \texttt{C5}): sits beside \texttt{U1}'s VREF pin.
\item \textbf{`converter`} (anchor \texttt{U1}; \texttt{C2}, \texttt{C3}, \texttt{C8}, \texttt{R7}): \texttt{C3} within 2.5 mm of
the VREF pin, SAME layer, no via between pad and cap - and what binds is the
charge LOOP (\texttt{C3} -\textgreater{} REF pin -\textgreater{} the internal array -\textgreater{} the converter's GND pin
-\textgreater{} back), so wide copper out AND back, and nothing in series between cap and
pin. \texttt{C2} within 2 mm of VDD with its ground via within 1 mm of its own pad -
that loop is the SPI drivers' current loop and must stay under 20 mm\textasciicircum{}2 (D3).
\texttt{R7} upstream of \texttt{C2}/\texttt{C8}, \texttt{C2} (the smaller) closest to the VDD pin.
\item \textbf{`entry`} (anchor \texttt{J2}; \texttt{C1}): the bulk cap at the rail entry, not at the
converter.
\end{itemize}
Spatial order across the board, from D4: {*}{*}\texttt{J1} -\textgreater{} \texttt{attenuator} -\textgreater{} \texttt{converter} + \texttt{reference} -\textgreater{} \texttt{J2}{*}{*}, with \texttt{J1} and \texttt{J2} on DIFFERENT edges, openings outward. Everything is on the FRONT side - single-sided assembly, and a part moved to the back would also cut the B.Cu return pour that D1 requires unbroken, so every ref carries an explicit \texttt{sides: front}.

\subsubsection*{6. Handover to P4}
1. Build the root sheet flat. Power symbols for \texttt{+3V3}, \texttt{GND}, \texttt{VREF} (the last one with its Value edited); root local labels for everything else. 2. \texttt{netlist\_audit.py --net ... --constraints ...} must come back clean before P5 touches the board. 3. Any net whose exported name differs from s2 is reconciled IN \texttt{constraints.json} at that moment, not later - an entry naming a net that does not exist fails silently, which is the whole reason this board is flat.

Full architecture narrative (not inlined; see the artifact index): \texttt{architecture/blocks.md}, \texttt{architecture/power\_tree.md}, \texttt{architecture/stackup.md}.
\section{Schematic}
\emph{Pending --- produced at P4.}
\section{Layout}
\emph{Pending --- produced at P6.}
\section{Verification}
\emph{Pending --- produced at P8.}
\section{DFM and Fabrication}
\emph{Pending --- produced at P9.}
\section{Run Record (Appendix)}
\subsection*{Phase digests}
\paragraph*{P0}
\begin{quote}\small\ttfamily\raggedright
\# P0 Intake - digest\\
~\\
- Mode `learning block-basics:` -\textgreater{} scope {*}{*}block-only{*}{*}, binding {*}{*}canonical{*}{*},\\
  stage none. Geometry is an OUTPUT: no size stated, none invented, `--outline auto`.\\
- `requirements.md` written; `check\_requirements.py` exit 0 / 0 violations, mode leg live.\\
- Owner answered all 9 questions in one batch: {*}{*}0.1 \% class{*}{*} at the terminal\\
  uncalibrated (+/-5 mV @25 C, +/-12 mV over 0-50 C), {*}{*}SPI{*}{*}, {*}{*}3.3 V only on J2{*}{*},\\
  low-voltage bench source, Q4/5/6/7/9 at defaults.\\
- The problem this board teaches, fixed at P0: 0.1 \% delivered THROUGH an attenuator\\
  that brings 0-5 V into a 3.3 V domain - ratio error and TCR are first-class terms\\
  beside the reference, not the converter's bit count.\\
- Safety: 5 V max on any net; mains/battery/motor/\textgreater{}30 V/high-I/RF do not apply, each\\
  contingent on the Q1 envelope (no isolation, no protection, input GND = host GND).\\
- Mode boundary for P1/P2: block-only excludes a second rail the block does not need\\
  {*}to work{*}; if 0-5 V from 3.3 V honestly needs one, that is support, and a decision.\\
- Owner then delegated design decisions ("decide yourself, best learning for the basic\\
  system"): H1-H4 become report-and-proceed; only a safety change re-blocks.\\
- Run-level: fable 5 out of credits -\textgreater{} every fable tier substitutes one step up.\\
- Defect: `modeslib.detect` skips `\textgreater{}`-lines, so a blockquoted brief silently disables\\
  the mode leg of the lint. Fixed here; `bb-mcu` has it too (reported, not touched).
\end{quote}
\paragraph*{P1}
\begin{quote}\small\ttfamily\raggedright
\# P1 Research - digest\\
~\\
Roster: 2x component-scout (`adc`, `afe-support`), reference-design (`afe`),\\
power-architect. No `research-interface-spec`: SPI-into-an-unspecified-host and a\\
screw terminal bind no standard. Detail lives in research/{*}.md.\\
~\\
- Topology = attenuator -\textgreater{} low-Ibias RRIO buffer -\textgreater{} SAR + external reference.\\
  Ratiometric-to-+3V3 dead (+/-5 \% rail = +/-250 mV vs a 5 mV budget, and nothing\\
  cancels: the source is not excited by that rail). Bufferless dead by\\
  t = (N+1){*}ln2{*}(Rs+Rss){*}Cs against the \textgreater{}=100 kohm input-impedance requirement.\\
- Leads: MCP3201 (C49274, separate VREF pin) / ADR4525A (C403698, 300 mV dropout\\
  leaves 335 mV at a 3.135 V rail) / LT5400B (C1739858) or a 0.01 \%-2 ppm discrete\\
  pair / OPA333 (C30878). MCP3202 disqualified (VDD and VREF share one pin in\\
  SOIC-8); every \textgreater{}=16-bit-ACCURACY SAR in stock needs \textgreater{}=4.5 V AVDD.\\
- Rail +3V3 `direct` from J2, 11 mA, 36 mW, no thermal constraints, no AVDD/DVDD\\
  ferrite (one net, 100 nF per pin, unsplit pour).\\
- Rulings recorded: RSS claim (4.8 mV) with the 10.6 mV worst-case published beside\\
  it; no second rail; Extended tier for all four accuracy parts; non-ratiometric;\\
  input range must include 0 V; ADI-branded only for AD8605-pattern parts.\\
- Open for P2: source impedance and divider Rtot are ONE number pair and the\\
  answered specs are jointly impossible at 0.1 \% (1 kohm into 100 kohm = 1 \%).\\
  Matched networks cap near 100 kohm/element - network-grade matching and\\
  megohm-grade input impedance cannot both be had.
\end{quote}
\paragraph*{P2}
\begin{quote}\small\ttfamily\raggedright
\# P2 Architecture + coverage research - digest\\
~\\
- Chain: screw terminal -\textgreater{} 5x 200k 0.02\%/10ppm equal string tapped 3:2 (K = 0.400,\\
  Rtot 1.00 Mohm) -\textgreater{} RRIO CMOS buffer -\textgreater{} ADS8326-class 16-bit SPI SAR, VREF 2.048 V\\
  from an ADR4520-class series reference. Full scale 5.12 V (2.4 \% headroom).\\
  Stackup JLC2313\_1.6, 2 layers, 1 oz. Flat root sheet. \textasciitilde{}\$30/board of parts.\\
- Budget: RSS 3.23 mV @25 C / 3.36 mV over 0-50 C against 5.0/12.0 targets;\\
  worst-case sum 8.35 / 10.29. Claim is made on RSS with the worst case published.\\
- Converter OVERRULED from P1: MCP3201's gain error is +/-5 LSB MAX = 6.10 mV at the\\
  terminal, 1.22x the whole budget. LSB-denominated DC specs shrink 16x from 12 to\\
  16 bits, so the higher-resolution part is the cheaper accuracy.\\
- Q9 NARROWED (visible owner item): source \textless{}= 200 ohm, board presents 1.00 Mohm, both\\
  on the silkscreen. As answered (\textless{}=1 kohm into \textgreater{}=100 kohm) the specs were jointly\\
  impossible at 0.1 \% - loading error is Rs/Rtot = 1 \% = 50 mV.\\
- The board's key architecture move: -IN wired as a remote sense to the string's\\
  bottom node. A divider passes ground offset at UNITY while dividing signal by K,\\
  so 1 mV of offset is 2.5 mV input-referred; the sense line cancels it exactly.\\
- Buffer: OPA333 fails at t\_acq = 9 us (GBW floor 48/(2{*}pi{*}t\_acq) = 849 kHz vs\\
  350 kHz; own settling 40.4 us). OPA320-class baseline; OPA333 revives if P3 finds\\
  acquisition is host-controlled. A buffer AT ALL is unconditional: switched-cap\\
  R\_eq = 1/(f\_s{*}C\_SH) \textasciitilde{} 2.08 Mohm vs a 240 kohm Thevenin is \textasciitilde{}10 \% gain error.\\
- Research: 4 tasks, 16 sources, {*}{*}32 verified records + 4 coverage checklists{*}{*} for\\
  topologies the library held nothing on. 3 correction rounds each; refutations that\\
  changed engineering: a reference topology read backwards off a layout figure, a\\
  stackup constant mistaken for a routing lever, an optimum treated as a maximum, an\\
  inverted curve ordering, a conversion window 5x off.\\
- Coverage after research: B2/B3/B4 provisional; B1 provisional except\\
  `constraints-emission`, recorded as a designed-under gap. Promotion DEFERRED to a\\
  separate owner pass - it writes the shared library and sibling runs are live.
\end{quote}
\subsection*{Gate attempt history}
\emph{(no entries)}
\subsection*{Issues}
\emph{(no entries)}
\subsection*{Human checkpoints}
{\small
\begin{longtable}{lllp{2.6cm}p{6.4cm}}
\toprule
\textbf{Id} & \textbf{Phase} & \textbf{Status} & \textbf{When} & \textbf{Note} \\
\midrule
\endhead
1 & P2 & not reached &  &  \\
2 & P4 & not reached &  &  \\
3 & P6 & not reached &  &  \\
4 & P8 & not reached &  &  \\
5 & P10 & not reached &  &  \\
\bottomrule
\end{longtable}
}
\section{Artifact Index}
{\small
\begin{longtable}{p{9cm}rp{4cm}}
\toprule
\textbf{File} & \textbf{Size} & \textbf{Note} \\
\midrule
\endhead
\texttt{state.json} & 79,162 B &  \\
\texttt{brief/brief.md} & 742 B &  \\
\texttt{requirements.md} & 23,802 B &  \\
\texttt{architecture/blocks.md} & 38,223 B &  \\
\texttt{architecture/constraints.json} & 27,161 B &  \\
\texttt{architecture/power\_tree.md} & 12,951 B &  \\
\texttt{architecture/sheets.md} & 14,973 B &  \\
\texttt{architecture/stackup.md} & 5,012 B &  \\
\bottomrule
\end{longtable}
}
\end{document}
